\section{DISCUSSÃO}
\subsection{O desvio pertence à comparação entre sistemas}

O resultado central é a natureza condicional da comparabilidade. A unidade de interpretação mais informativa não é uma tecnologia isolada, mas a combinação entre instrumento, configuração, comparador, sítio e período. Os contrastes de TEOM convencional, SES, sistema diferencial e FDMS mostram como mudanças no tratamento da amostra alteram a relação com a gravimetria \parencite{rsl201,rsl144,rsl101,rsl105}. Agrupar todas essas configurações sob um único nome pode fazer resultados distintos parecerem contraditórios quando, na realidade, medem depósitos formados e preservados de maneira diferente.

Essa interpretação exige simetria na apreciação dos métodos. Um automático abaixo do filtro pode ter perdido semivoláteis; um automático acima do filtro pode estar preservando material que o comparador perdeu. Água retida ou adsorvida e diferenças de seleção por tamanho também podem atuar nos dois lados da comparação. Em Rubidoux, a evidência química de nitrato é importante justamente porque vai além de uma discrepância numérica; em Erfurt, o baixo nitrato aponta para outra composição semivolátil \parencite{rsl101,rsl201}. O mecanismo precisa ser confrontado com as condições observadas, sem assumir que todo desvio do mesmo sinal tem a mesma causa.

O desenho define quais componentes dessa cadeia foram examinados. A comparação beta com a pesagem do mesmo filtro reduz diferenças de coleta e é adequada para investigar a determinação da massa, mas deixa de testar parte da cadeia de um sistema independente \parencite{rsl214}. Na avaliação dos pDR, a distinção entre massa coletada, resposta óptica e gravimetria de referência oferece informação que se perderia em um único coeficiente global \parencite{rsl089}. Esses desenhos são complementares, desde que seus resultados não sejam tratados como ensaios equivalentes.

Também não há base para um ranking único entre frações ou princípios. Os melhores resultados de SMPS--APS na fração fina não resolvem a seleção e as hipóteses de massa da fração grossa; a maior média de desvio dos monitores ópticos na análise de rede norte-americana não estima o desempenho de todo fotômetro em qualquer aerossol \parencite{rsl288,rsl283,rsl026}. As campanhas de Como, Riohacha e Baltimore mostram que a própria direção da diferença óptica muda conforme equipamento e contexto \parencite{rsl388,rsl226,rsl368}.

\subsection{Indicadores, correções e domínio de transferência}

Uma parte importante das divergências aparentes decorre da definição da estatística. Com concentrações gravimétricas positivas, a razão entre médias dá maior influência aos pares com maior concentração gravimétrica; a média das razões atribui o mesmo peso a cada razão e se torna particularmente sensível a denominadores pequenos. A média geométrica resume as razões na escala logarítmica. Nenhuma dessas medidas é intercambiável com inclinação de regressão, desvio-padrão ou incerteza de medição. Por isso, aproximar numericamente percentuais publicados por estudos diferentes pode criar uma semelhança que não existe na definição do resultado.

O contraste entre o TEOM de Bakersfield, com alta correlação e inclinação distante da unidade, e os ensaios de massa no mesmo filtro demonstra a necessidade de ler associação, escala e desenho em conjunto \parencite{rsl129,rsl214}. Um coeficiente de determinação elevado descreve variação explicada pela relação ajustada; não estabelece, isoladamente, que diferenças sejam pequenas para a finalidade pretendida \parencite{bland1986}. Em baixas concentrações, o intercepto e a dispersão residual ganham importância mesmo quando a inclinação se aproxima de um.

As correções acrescentam uma distinção entre reproduzir um comparador e conservar as características do aerossol originalmente amostrado. A correção de um sistema que retém mais semivoláteis para aproximá-lo de um filtro sujeito à volatilização pode melhorar a concordância entre séries sem demonstrar maior proximidade de uma massa ambiental sem artefatos \parencite{rsl101,rsl085}. A referência e a finalidade da correção devem, portanto, ser explícitas. No estudo de \textcite{rsl085}, a própria expressão \enquote{massa real}, associada ao FDMS pelos autores, corresponde à interpretação operacional adotada no estudo e não foi assumida nesta revisão como verdade independente do sistema de medição.

Os ajustes podem incorporar temperatura, umidade, concentração ou resposta volátil, e modelos mais flexíveis podem representar relações que um fator constante não captura \parencite{rsl024,rsl029,rsl061,rsl219}. Contudo, aumentar $R^2$ ou reduzir uma média de desvio no conjunto usado para desenvolver o ajuste não estabelece desempenho futuro. No procedimento de \textcite{rsl085}, a aplicação ao conjunto total inclui pontos utilizados na derivação. Em outros relatos, o uso de calibração de um sítio em outro fornece uma situação de transferência a examinar, sem estabelecer aplicabilidade irrestrita \parencite{rsl342}. A conclusão apropriada deve acompanhar o período, o sítio e as condições efetivamente testados.

A transparência sobre a seleção dos pares também é necessária. As quatro exclusões antes da regressão SHARP--FRM e a diferença entre os 20 pares brutos e os 18 corrigidos na comparação óptica não invalidam os resultados; delimitam sua interpretação \parencite{rsl059,rsl089}. Uma avaliação de correção deve apresentar o desempenho dos conjuntos comparados, as regras de exclusão e a dispersão remanescente, para distinguir ganho de ajuste de mudanças na composição da amostra.

\subsection{Integração de séries e representatividade temporal}

Para integrar dados de redes ou períodos, a comparabilidade deve ser tratada antes da atribuição das diferenças à evolução ambiental. Mudança de tecnologia, temperatura de operação, secagem ou algoritmo pode alterar a distribuição das concentrações e introduzir uma ruptura instrumental na série. Uma diferença que varia com estação, composição ou concentração não é necessariamente removida por um único fator aplicado ao ano inteiro. Os resultados de campo e a heterogeneidade entre monitores da rede norte-americana sustentam essa preocupação \parencite{rsl049,rsl144,rsl026}.

A agregação temporal também precisa preservar a base de medição. Uma média de valores horários deve corresponder ao intervalo em que o filtro efetivamente coletou amostra; diferenças de cobertura, relógio e seleção dos dados podem modificar os pares avaliados. A interpretação deve distinguir comparações diárias, integrações de vários dias e médias anuais. Passar de uma escala para outra altera a contribuição de episódios, a dispersão e o significado do resultado; não acrescenta, por si, evidência de concordância no intervalo original \parencite{rsl129,rsl188,rsl026}.

Essa questão afeta diretamente análises de frequência e distribuição temporal de amostragem. Se o desvio depende da meteorologia ou do tipo de episódio, selecionar dias ou horários também seleciona condições em que a relação instrumental pode mudar. Uma diferença atribuída à menor frequência de coleta pode conter componente associado à mudança de método ou à aplicação de correções em condições diferentes da calibração. Assim, o exame da representatividade temporal requer documentar simultaneamente o calendário amostral e a continuidade do sistema de medição.

O Quadro~\ref{qua:aplicacao} traduz os fatores encontrados em informações que permitem interpretar e utilizar comparações. Essa organização não fornece um critério novo de equivalência. Ela conecta mecanismos e desenhos observados a decisões sobre agregação, correção e integração de séries. A aplicação a redes locais, inclusive em Pernambuco, depende da identificação de seus equipamentos e das condições meteorológicas e de aerossol efetivamente representadas; nenhuma equação foi validada para o estado por esta revisão.

\begin{quadro*}[tbp]
\caption{Condições necessárias para interpretar e utilizar comparações entre métodos}\label{qua:aplicacao}
\begingroup\singlespacing\fontsize{10}{12}\selectfont
\begin{tabularx}{\textwidth}{|P{.23\textwidth}|P{.34\textwidth}|X|}
\hline\textbf{Condicionante} & \textbf{Informação a preservar} & \textbf{Implicação para uso dos dados}\\\hline
Aquecimento e semivoláteis & Temperaturas, secagem e retenção ou compensação de material nos dois sistemas \parencite{rsl201,rsl101,rsl144}. & Interpretar a diferença em relação ao tratamento da amostra, incluindo os artefatos do filtro.\\\hline
Água e propriedades ópticas & Umidade, composição, configuração e identificação das séries bruta e corrigida \parencite{rsl170,rsl089,rsl388}. & Delimitar em quais condições uma correção de umidade melhora a relação, sem assumir resposta única por tecnologia.\\\hline
Seleção por tamanho e conversão em massa & Fração, entrada, densidade e forma assumidas na conversão \parencite{rsl136,rsl288,rsl283}. & Evitar transportar desempenho de uma fração ou comparador para outro sem avaliação.\\\hline
Concentração e definição da estatística & Diferenças absolutas e relativas, intercepto, dispersão e tipo de razão \parencite{rsl073,rsl129,rsl085}. & Examinar baixas concentrações e distinguir associação de concordância.\\\hline
Pareamento e escala temporal & Intervalos efetivamente coincidentes, cobertura, exclusões e denominador por contraste \parencite{rsl059,rsl026}. & Separar comportamento no intervalo de coleta e resultado agregado; registrar rupturas instrumentais na série.\\\hline
Desenvolvimento e transferência de correções & Dados de desenvolvimento, dados de avaliação, forma do modelo e domínio ambiental \parencite{rsl029,rsl085,rsl219,rsl342}. & Restringir a aplicação ao alcance demonstrado e examinar o desempenho em condições distintas do ajuste.\\\hline
\end{tabularx}
\Fonte{Elaboração própria a partir dos condicionantes e desenhos discutidos nos relatos citados.}
\NotaQuadro{O quadro organiza a interpretação e não estabelece um protocolo novo de equivalência ou uma classificação de risco dos estudos.}
\endgroup
\end{quadro*}


\subsection{Limitações da evidência e prioridades de pesquisa}

A cobertura executada restringe o alcance do corpus. As três bases, a busca delimitada de citações, os 32 textos não recuperados e as seis identificações insuficientes podem ter deixado de fora evidências relevantes. O encerramento e as delimitações metodológicas ocorreram após conhecimento dos resultados. A concentração de relatos em determinados períodos de publicação descreve o conjunto selecionado e não autoriza concluir que a pesquisa sobre comparabilidade diminuiu ou que uma tecnologia recente tenha pouca evidência em toda a literatura.

A extração focal preservou estimativas pertinentes, mas não desdobrou todas as combinações de local, período, instrumento e comparador. Relatos curtos e apresentações apenas agregadas limitaram algumas interpretações. A apreciação narrativa registrou limitações de desenho, transparência e participação de fabricantes, sem converter uma declaração de financiamento em prova de viés nem atribuir categorias uniformes de risco ao conjunto ampliado. A dependência potencial entre publicações também restringe o uso das contagens como tamanho efetivo de uma base de experimentos.

A condução por pesquisador único com assistência de inteligência artificial mantém suscetibilidade a erros de leitura, extração e julgamento. A rastreabilidade das estimativas, a distinção entre resultados próprios e citações bibliográficas e a explicitação dos denominadores tornam essas decisões verificáveis, mas não substituem uma revisão independente. Essas limitações, somadas à heterogeneidade dos desenhos e indicadores, restringem a precisão das conclusões e sustentam a opção por interpretação narrativa, sem classificação global de desempenho.

Os avanços mais úteis são comparações que documentem entradas e tratamentos dos dois sistemas, caracterizem componentes relevantes do aerossol e preservem os intervalos e denominadores de cada resultado. A apresentação conjunta de diferenças absolutas e relativas, dispersão e critérios de concordância permitiria avaliar melhor concentrações baixas e episódios. Para correções, períodos ou sítios efetivamente separados do desenvolvimento ajudariam a delimitar transferência. A contribuição esperada é reduzir a ambiguidade do contraste e identificar quando uma diferença instrumental pode comprometer a interpretação ambiental.
